% --- CORRECCIÓN IMPORTANTE EN LA PRIMERA LÍNEA ---
% Agregamos "xcolor=table" para que funcione \rowcolor sin errores
\documentclass[aspectratio=169, 10pt, xcolor=table]{beamer}

% --- Configuración del Tema (estilo profesional como el ejemplo) ---
\usetheme{Madrid}
\usecolortheme{dolphin}

% --- Paquetes necesarios ---
\usepackage[utf8]{inputenc}
\usepackage[spanish]{babel}
\usepackage{amsmath, amssymb}
\usepackage{booktabs}
\usepackage{graphicx}
\usepackage{listings}
\usepackage{tikz}
\usetikzlibrary{arrows.meta, positioning, shapes.geometric}

% --- Ruta de Imágenes (crea la carpeta "imagenes" y coloca ahí tus archivos) ---
\graphicspath{{./imagenes/}}

% --- Estilo para código Python ---
\definecolor{codegreen}{rgb}{0,0.6,0}
\definecolor{codegray}{rgb}{0.5,0.5,0.5}
\definecolor{codepurple}{rgb}{0.58,0,0.82}
\definecolor{backcolour}{rgb}{0.95,0.95,0.92}
\lstset{
    language=Python,
    backgroundcolor=\color{backcolour},
    commentstyle=\color{codegreen},
    keywordstyle=\color{blue},
    numberstyle=\tiny\color{codegray},
    stringstyle=\color{codepurple},
    basicstyle=\ttfamily\scriptsize,
    breaklines=true,
    frame=single,
    showstringspaces=false
}

% --- Pie de página personalizado ---
\makeatletter

\setbeamertemplate{footline}{
  \leavevmode%
  \hbox{%
  \begin{beamercolorbox}[wd=.333333\paperwidth,ht=2.25ex,dp=1ex,center]{author in head/foot}%
    \usebeamerfont{author in head/foot}\insertshortauthor
  \end{beamercolorbox}%
  \begin{beamercolorbox}[wd=.333333\paperwidth,ht=2.25ex,dp=1ex,center]{title in head/foot}%
    \usebeamerfont{title in head/foot}Sesión 11
  \end{beamercolorbox}%
  \begin{beamercolorbox}[wd=.333333\paperwidth,ht=2.25ex,dp=1ex,right]{date in head/foot}%
    \usebeamerfont{date in head/foot}Machine Learning \hfill \insertframenumber/\inserttotalframenumber \hspace*{0.3cm}
  \end{beamercolorbox}}%
  \vskip0pt%
}

\makeatother

% --- Datos de la portada ---
\title[SESIÓN 11]{\textbf{Sesión 11}}
\subtitle{Modelos Complementarios: \\ Técnicas especializadas para problemas específicos}
\author[Prof. C. Sánchez]{Carlos César Sánchez Coronel}
\institute{\textbf{MACHINE LEARNING}}
\date{}

% Logo opcional (descomentar si tienes la imagen)
% \titlegraphic{\includegraphics[height=1.5cm]{logo.png}}

% --- Eliminar la barra de navegación (opcional) ---
\setbeamertemplate{navigation symbols}{}

% --- Índice al inicio de cada sección ---
\AtBeginSection[]{
  \begin{frame}
    \frametitle{Contenido}
    \tableofcontents[currentsection]
  \end{frame}
}

% ============================================================================
% INICIO DEL DOCUMENTO
% ============================================================================
\begin{document}

% --- Portada ---
\begin{frame}
    \titlepage
\end{frame}

% --- Índice general ---
\begin{frame}{Contenido}
    \tableofcontents
\end{frame}

% ============================================================================
% SECCIÓN 1: REGRESIÓN POLINÓMICA
% ============================================================================
\section{Regresión Polinómica}

\begin{frame}{Regresión polinómica}
    \begin{block}{Idea}
        Añadir términos polinómicos de las variables originales como nuevas características para capturar relaciones no lineales.
    \end{block}
    Para una variable $x$, grado $d$:
    \[
    y = \beta_0 + \beta_1 x + \beta_2 x^2 + \cdots + \beta_d x^d + \varepsilon
    \]
    Sigue siendo un modelo \textbf{lineal en los parámetros} $\beta$ → se estima por mínimos cuadrados.
    % Basado en Pregunta 1
    % IMAGEN: Ajuste polinómico de diferentes grados
\end{frame}

\begin{frame}{Riesgo de sobreajuste}
    \begin{itemize}
        \item A mayor grado $d$, mayor flexibilidad.
        \item Riesgo de \textbf{overfitting}, especialmente en los extremos.
        \item Se elige el grado mediante validación cruzada.
    \end{itemize}
    \begin{exampleblock}{Ejemplo (Pregunta 2)}
        RMSE en validación: grado1=4.2, grado2=3.1, grado3=2.8, grado4=2.7, grado5=3.5.\\
        Se elige grado 3 o 4 (trade-off sesgo-varianza).
    \end{exampleblock}
    \begin{itemize}
        \item Regularización (Ridge, Lasso) permite grados altos sin sobreajuste.
    \end{itemize}
\end{frame}

% ============================================================================
% SECCIÓN 2: REGRESIÓN ROBUSTA (HUBER, RANSAC)
% ============================================================================
\section{Regresión Robusta}

\begin{frame}{Regresión de Huber}
    \begin{block}{Función de pérdida}
        \[
        L_{\delta}(r) = 
        \begin{cases}
            \frac{1}{2}r^2 & \text{si } |r| \le \delta \\
            \delta(|r| - \frac{1}{2}\delta) & \text{en otro caso}
        \end{cases}
        \]
        donde $r = y - \hat{y}$.
    \end{block}
    \begin{itemize}
        \item Combina MSE (para errores pequeños) y MAE (para errores grandes).
        \item Menos sensible a outliers que MSE, más eficiente que MAE.
        \item Parámetro $\delta$ controla la transición (típico $\delta = 1.35$ da 95\% eficiencia si los errores son normales).
    \end{itemize}
    % Basado en Pregunta 3
\end{frame}

\begin{frame}{RANSAC (RANdom SAmple Consensus)}
    \begin{block}{Algoritmo}
        \begin{enumerate}
            \item Seleccionar aleatoriamente una muestra mínima para ajustar un modelo.
            \item Calcular cuántos puntos están dentro de un umbral de distancia (inliers).
            \item Repetir $N$ iteraciones, guardar el modelo con más inliers.
            \item Reajustar el modelo final con todos los inliers.
        \end{enumerate}
    \end{block}
    \begin{itemize}
        \item Extremadamente robusto (tolera hasta >50\% de outliers).
        \item Parámetros: tamaño de muestra mínimo, umbral de distancia, número de iteraciones.
    \end{itemize}
    % Basado en Pregunta 4
\end{frame}

\begin{frame}{Comparación de pérdidas}
    \begin{table}
        \centering
        \begin{tabular}{l|c|c}
            \toprule
            \textbf{Función} & \textbf{Ecuación} & \textbf{Sensibilidad a outliers} \\
            \midrule
            MSE & $r^2$ & Alta \\
            MAE & $|r|$ & Moderada \\
            Huber & $L_{\delta}(r)$ & Baja (con $\delta$ adecuado) \\
            \bottomrule
        \end{tabular}
    \end{table}
    % Basado en Pregunta 5
\end{frame}

% ============================================================================
% SECCIÓN 3: SUPPORT VECTOR REGRESSION (SVR)
% ============================================================================
\section{Support Vector Regression (SVR)}

\begin{frame}{SVR: formulación}
    \begin{block}{Objetivo}
        Encontrar una función $f(x) = w^T x + b$ que tenga una desviación máxima $\varepsilon$ respecto a los valores reales, y sea lo más plana posible.
    \end{block}
    \begin{columns}[t]
        \column{0.5\textwidth}
        \textbf{Primal}
        \[
        \min \frac{1}{2}\|w\|^2 + C\sum (\xi_i + \xi_i^*)
        \]
        \column{0.5\textwidth}
        \textbf{Sujeto a}
        \begin{align*}
            y_i - f(x_i) &\le \varepsilon + \xi_i \\
            f(x_i) - y_i &\le \varepsilon + \xi_i^* \\
            \xi_i, \xi_i^* &\ge 0
        \end{align*}
    \end{columns}
    % Basado en Pregunta 6
\end{frame}

\begin{frame}{Interpretación de parámetros}
    \begin{itemize}
        \item $\varepsilon$: ancho de la banda de tolerancia (errores menores no penalizan).
        \item $C$: regularización; controla penalización por puntos fuera de la banda.
        \item Kernel: permite relaciones no lineales (RBF, polinomial, etc.).
    \end{itemize}
    \begin{exampleblock}{SVR con kernel RBF}
        \begin{itemize}
            \item $\gamma$ (gamma) controla la influencia de cada punto.
            \item Valores grandes de $\gamma$ pueden causar sobreajuste.
            \item Se ajustan con validación cruzada.
        \end{itemize}
    \end{exampleblock}
    % Basado en Pregunta 7
\end{frame}

% ============================================================================
% SECCIÓN 4: DBSCAN A FONDO
% ============================================================================
\section{DBSCAN a fondo}

\begin{frame}{DBSCAN: elección de parámetros}
    \begin{itemize}
        \item \textbf{minPts}: regla heurística $\ge$ dimensionalidad + 1 (típico 3-10).
        \item \textbf{eps ($\varepsilon$)}: se elige mediante el gráfico de distancia al k-ésimo vecino (k = minPts).
        \begin{enumerate}
            \item Calcular distancia al k-ésimo vecino para cada punto.
            \item Ordenar y graficar.
            \item Elegir eps en el ``codo'' de la curva.
        \end{enumerate}
    \end{itemize}
    % Basado en Pregunta 9
    % IMAGEN: Gráfico de distancias al k-ésimo vecino
\end{frame}

\begin{frame}{Variantes de DBSCAN}
    \begin{description}
        \item[OPTICS] No requiere un valor fijo de $\varepsilon$; produce un diagrama de alcanzabilidad que revela estructuras de densidad variable.
        \item[HDBSCAN] Extensión jerárquica que encuentra clusters de diferentes densidades y es más robusta a los parámetros.
    \end{description}
    % Basado en Pregunta 10
\end{frame}

% ============================================================================
% SECCIÓN 5: CLUSTERING JERÁRQUICO A FONDO
% ============================================================================
\section{Clustering Jerárquico a fondo}

\begin{frame}{Tipos de enlace}
    \begin{columns}[t]
        \column{0.5\textwidth}
        \textbf{Single linkage}
        \begin{itemize}
            \item Distancia mínima entre clusters.
            \item Tiende a formar cadenas (efecto ``encadenamiento'').
        \end{itemize}
        \textbf{Complete linkage}
        \begin{itemize}
            \item Distancia máxima.
            \item Clusters compactos y esféricos.
        \end{itemize}
        \column{0.5\textwidth}
        \textbf{Average linkage}
        \begin{itemize}
            \item Distancia promedio.
            \item Compromiso.
        \end{itemize}
        \textbf{Ward linkage}
        \begin{itemize}
            \item Minimiza incremento de varianza intra-cluster.
            \item Clusters de tamaño similar.
        \end{itemize}
    \end{columns}
    % Basado en Pregunta 11
    % IMAGEN: Dendrograma con diferentes enlaces
\end{frame}

\begin{frame}{Corte del dendrograma}
    \begin{itemize}
        \item Se corta a una altura donde las distancias de fusión son grandes.
        \item La elección puede basarse en:
        \begin{itemize}
            \item Criterio visual (saltos grandes).
            \item Silueta promedio para diferentes cortes.
        \end{itemize}
        \item Complejidad computacional: $O(n^3)$ (ingenuo) o $O(n^2 \log n)$ con optimizaciones.
    \end{itemize}
    % Basado en Pregunta 12
\end{frame}

% ============================================================================
% SECCIÓN 6: MÉTRICAS ADICIONALES PARA CLUSTERING
% ============================================================================
\section{Métricas adicionales para clustering}

\begin{frame}{Índice de Rand y ARI}
    \begin{itemize}
        \item \textbf{Rand Index (RI)}: mide concordancia entre dos particiones.
        \[
        RI = \frac{a+b}{\binom{n}{2}}
        \]
        donde $a$ = pares en mismo cluster en ambas, $b$ = pares en diferentes clusters en ambas.
        \item \textbf{Adjusted Rand Index (ARI)}: corrige por el azar, dando 0 para particiones aleatorias, 1 para coincidencia perfecta.
    \end{itemize}
    % Basado en Pregunta 13
\end{frame}

\begin{frame}{Información Mutua (MI) y NMI}
    \begin{itemize}
        \item \textbf{MI}: mide la información compartida entre dos particiones.
        \[
        MI(U,V) = \sum_{i,j} P(i,j) \log \frac{P(i,j)}{P(i)P(j)}
        \]
        \item \textbf{NMI (Normalized MI)}: escala a $[0,1]$ dividiendo por la media de las entropías.
    \end{itemize}
    % Basado en Pregunta 14
\end{frame}

\begin{frame}{Comparación de métricas de clustering}
    \begin{table}
        \centering
        \begin{tabular}{l|c|c|c}
            \toprule
            \textbf{Métrica} & \textbf{Rango} & \textbf{Corrige azar} & \textbf{Requiere etiquetas} \\
            \midrule
            Silueta & $[-1,1]$ & No & No \\
            Davies-Bouldin & $[0,\infty)$ (menor mejor) & No & No \\
            Rand Index & $[0,1]$ & No & Sí \\
            ARI & $[-1,1]$ & Sí & Sí \\
            NMI & $[0,1]$ & Aproximadamente & Sí \\
            \bottomrule
        \end{tabular}
    \end{table}
    % Basado en Pregunta 15
\end{frame}

% ============================================================================
% SECCIÓN 7: GUÍA DE SELECCIÓN
% ============================================================================
\section{Guía de selección de técnicas}

\begin{frame}{Cuándo usar cada técnica}
    \begin{table}
        \centering
        \begin{tabular}{p{4cm}|p{6cm}}
            \toprule
            \textbf{Técnica} & \textbf{Cuándo usarla} \\
            \midrule
            Regresión polinómica & Relación no lineal suave, pocas variables, interpretabilidad. \\
            Regresión robusta (Huber) & Outliers moderados, se desea eficiencia. \\
            RANSAC & Gran cantidad de outliers (hasta >50\%), modelo de inliers. \\
            SVR & Relaciones no lineales, robustez a outliers, control sobre banda de error. \\
            DBSCAN & Clusters de forma arbitraria, presencia de ruido, no se conoce $k$. \\
            Clustering jerárquico & Se necesita jerarquía interpretable, pocos datos. \\
            \bottomrule
        \end{tabular}
    \end{table}
    % Basado en teoría
\end{frame}

% ============================================================================
% SECCIÓN 8: CASO INTEGRADO - SENSORES INDUSTRIALES
% ============================================================================
\section{Caso integrado: Sensores industriales}

\begin{frame}{Problema}
    \begin{itemize}
        \item Planta industrial con sensores de temperatura, presión, vibración (10,000 lecturas, sin etiquetas).
        \item Objetivos:
        \begin{enumerate}
            \item Identificar modos de operación normal (clusters).
            \item Detectar anomalías (posibles fallos incipientes).
            \item Modelar relaciones esperadas entre sensores.
        \end{enumerate}
    \end{itemize}
    % Basado en Pregunta 20
\end{frame}

\begin{frame}{Pipeline propuesto}
    \begin{enumerate}
        \item \textbf{DBSCAN} para encontrar clusters de operación normal (puntos densos). El ruido son anomalías candidatas.
        \item Para cada cluster, ajustar una \textbf{regresión robusta} (Huber o RANSAC) modelando, ej., temperatura vs presión.
        \item En tiempo real:
        \begin{itemize}
            \item Si un punto no pertenece a ningún cluster normal $\rightarrow$ alerta.
            \item Si el residuo de la regresión supera un umbral $\rightarrow$ alerta.
        \end{itemize}
    \end{enumerate}
    % Basado en Pregunta 16
\end{frame}

\begin{frame}{Elección de parámetros y validación}
    \begin{itemize}
        \item DBSCAN: elegir minPts (ej. 5) y eps con gráfico de distancia al 5-vecino.
        \item Para clusters normales, regresión robusta (Huber) o polinómica si la relación es no lineal.
        \item Si se tienen pocos eventos de fallo conocidos, usar \textbf{ARI/NMI} para validar que los clusters capturan estados pre-falla.
    \end{itemize}
\end{frame}

% ============================================================================
% SECCIÓN 9: IMPLEMENTACIÓN EN PYTHON
% ============================================================================
\section{Implementación en Python}

\begin{frame}[fragile]{Regresión polinómica y robusta}
    \begin{lstlisting}
from sklearn.preprocessing import PolynomialFeatures
from sklearn.linear_model import LinearRegression, HuberRegressor, RANSACRegressor
from sklearn.pipeline import Pipeline

# Polinómica
model = Pipeline([
    ('poly', PolynomialFeatures(degree=3)),
    ('linear', LinearRegression())
])
model.fit(X, y)

# Huber
huber = HuberRegressor(epsilon=1.35)
huber.fit(X, y)

# RANSAC
ransac = RANSACRegressor(min_samples=50, residual_threshold=5.0, max_trials=100)
ransac.fit(X, y)
inlier_mask = ransac.inlier_mask_
    \end{lstlisting}
\end{frame}

\begin{frame}[fragile]{SVR y DBSCAN}
    \begin{lstlisting}
from sklearn.svm import SVR
from sklearn.cluster import DBSCAN
from sklearn.neighbors import NearestNeighbors

# SVR
svr = SVR(kernel='rbf', C=1.0, epsilon=0.1)
svr.fit(X_train, y_train)

# Gráfico de distancias para DBSCAN
neigh = NearestNeighbors(n_neighbors=5)
neigh.fit(X)
distances, _ = neigh.kneighbors(X)
k_dist = np.sort(distances[:, -1])
plt.plot(k_dist); plt.show()

# DBSCAN
dbscan = DBSCAN(eps=0.5, min_samples=5)
labels = dbscan.fit_predict(X)
    \end{lstlisting}
\end{frame}

\begin{frame}[fragile]{Clustering jerárquico y métricas externas}
    \begin{lstlisting}
from scipy.cluster.hierarchy import dendrogram, linkage, fcluster
from sklearn.metrics import adjusted_rand_score, normalized_mutual_info_score

# Jerárquico
Z = linkage(X, method='ward')
plt.figure(figsize=(10,5))
dendrogram(Z)
plt.show()

clusters = fcluster(Z, t=3, criterion='maxclust')

# Métricas externas (con etiquetas verdaderas y_true)
ari = adjusted_rand_score(y_true, clusters)
nmi = normalized_mutual_info_score(y_true, clusters)
    \end{lstlisting}
\end{frame}

% ============================================================================
% SECCIÓN 10: CONEXIÓN Y RESUMEN
% ============================================================================
\section{Conexión y Resumen}

\begin{frame}{Conexión con el resto del curso}
    \begin{itemize}
        \item Regresión robusta y SVR complementan los modelos de regresión.
        \item DBSCAN y clustering jerárquico amplían opciones de clustering (K-Means).
        \item Métricas externas (ARI, NMI) permiten evaluar clustering con ground truth.
    \end{itemize}
\end{frame}

\begin{frame}{Lo que hemos aprendido}
    \begin{exampleblock}{Conceptos clave}
        \begin{itemize}
            \item Regresión polinómica: no linealidades, sobreajuste, regularización.
            \item Regresión robusta: Huber, RANSAC, comparación de pérdidas.
            \item SVR: banda $\varepsilon$, parámetros $C$, $\gamma$, kernel.
            \item DBSCAN: elección de $\varepsilon$ y minPts, variantes OPTICS/HDBSCAN.
            \item Clustering jerárquico: tipos de enlace, dendrograma, corte.
            \item Métricas externas: ARI, NMI.
            \item Caso práctico: detección de anomalías en sensores industriales.
        \end{itemize}
    \end{exampleblock}
\end{frame}

% --- Diapositiva final ---
\begin{frame}
    \centering
    \Huge \textbf{¡Gracias!}

\end{frame}

\end{document}